\documentclass[UTF8,11pt,a4paper]{ctexart}
\usepackage{amsmath,amssymb}
\usepackage{booktabs,longtable,array}
\usepackage[a4paper,landscape,left=1.4cm,right=1.4cm,top=1.5cm,bottom=1.5cm]{geometry}
\usepackage[hidelinks]{hyperref}
\usepackage{fancyhdr}
\newcommand{\keywords}[1]{\par\noindent\textbf{关键词：}#1}
\setlength{\headheight}{14pt}
\setlength{\parindent}{2em}
\setlength{\parskip}{0.25em}
\setlength{\LTpre}{0.35em}
\setlength{\LTpost}{0.55em}
\setlength{\tabcolsep}{3pt}
\renewcommand{\arraystretch}{1.22}
\pagestyle{fancy}
\fancyhf{}
\fancyhead[L]{2024 年高教社杯 C 题：农作物种植策略}
\fancyhead[R]{问题 1--3 总思路表}
\fancyfoot[C]{\thepage}
\setcounter{secnumdepth}{0}
\begin{document}
\title{C题总思路表}
\author{}
\date{}
\maketitle
\begin{abstract}
本文将农作物种植策略问题一至问题三的数据处理、数学模型、求解流程、验证方法、现有结果及证据边界汇总为统一思路表，供后续工程实现与论文写作核对使用。
\keywords{农作物种植；混合整数线性规划；情景分析；CVaR；t-Copula}
\end{abstract}

\begin{quote}\small
用途：作为论文写作底稿，记录问题 1—3 最终采用的数据处理、模型、求解、验证和结论口径。数值结果只引用当前工程中的真实运行记录；问题 3 尚未编程求解，不填写虚构结果。
\end{quote}

\section{一、总体路线}

\noindent 本节汇总见表~\ref{tab:thought-1}。
\begin{longtable}{@{}>{\raggedright\arraybackslash}p{0.20\linewidth}>{\raggedright\arraybackslash}p{0.36\linewidth}>{\raggedright\arraybackslash}p{0.38\linewidth}@{}}
\caption{思路表第 1 表}\label{tab:thought-1}\\
\toprule
\textbf{阶段} & \textbf{核心任务} & \textbf{当前状态} \\
\midrule
\endfirsthead
\toprule
\textbf{阶段} & \textbf{核心任务} & \textbf{当前状态} \\
\midrule
\endhead
1. 数据清洗 & 统一附件字段、补齐合并单元格、建立地块—作物—季次适种关系，并由 2023 年产量构造销量代理 & 已完成 \\
2. 确定性可行域 & 建立面积、适种、水浇地模式、重茬、滚动三年豆类和最小面积约束 & 已完成 \\
3. 问题 1 确定性优化 & 分别求解超产浪费与超产半价出售两种 MILP & 已运行；有可行解，均未完成最优性认证 \\
4. 问题 2 随机优化 & LHS 情景生成、尾部保护缩减、均值—下尾 CVaR 和三级字典序求解 & 已运行；前沿不完整，选中方案未认证 \\
5. 问题 3 相关与交互扩展 & t-Copula 相关情景、交叉价格弹性、豆类前茬互补和联合随机 MILP & 建模方案已完成，等待编程求解 \\
6. 验证与稳健性 & 约束复算、利润复算、样本外评估、相关性审计、消融与参数敏感性 & 问题 1、2 部分完成；问题 3 待实现 \\
7. 最终建议 & 汇总三问方案、最优性边界、适用条件和局限 & 待问题 3 数值结果完成后冻结 \\
\bottomrule
\end{longtable}

\section{二、数据清洗}

\noindent 本节汇总见表~\ref{tab:thought-2}。
\begin{longtable}{@{}>{\raggedright\arraybackslash}p{0.05\linewidth}>{\raggedright\arraybackslash}p{0.15\linewidth}>{\raggedright\arraybackslash}p{0.27\linewidth}>{\raggedright\arraybackslash}p{0.16\linewidth}>{\raggedright\arraybackslash}p{0.25\linewidth}>{\raggedright\arraybackslash}p{0.08\linewidth}@{}}
\caption{思路表第 2 表}\label{tab:thought-2}\\
\toprule
\textbf{编号} & \textbf{数据问题} & \textbf{当前处理方案} & \textbf{处理对象/阶段} & \textbf{处理理由} & \textbf{状态} \\
\midrule
\endfirsthead
\toprule
\textbf{编号} & \textbf{数据问题} & \textbf{当前处理方案} & \textbf{处理对象/阶段} & \textbf{处理理由} & \textbf{状态} \\
\midrule
\endhead
1 & 附件中存在合并单元格、空白地块名和分组字段 & 只读加载后对合并区域产生的空白地块名与作物分组向下填充，再用编号和适种说明交叉核验 & 附件 1、附件 2 & 避免把展示性空白误认为数据缺失 & 已确认 \\
2 & “菠菜 ”、“生菜 ”、“普通大棚 ”等文本带尾随空格，且出现 \texttt{\_x000d\_} & 统一去除首尾空格，将 \texttt{\_x000d\_} 还原为空格或换行，并以规范化名称建立键 & 作物名、地块类型、季次 & 防止同一对象被拆成多个类别或适种匹配失败 & 已确认 \\
3 & 面积、编号、亩产量、成本、价格端点可能以数值字符串保存 & 显式转换为数值；无法转换、非有限或超出合理范围时停止运行 & 三问全部参数 & 不依赖 Excel 显示格式猜测类型 & 已确认 \\
4 & 智慧大棚第一季统计参数在附件 2 中省略 & 按附件 2 说明继承普通大棚第一季同种蔬菜的亩产量、成本和价格；同时补齐 2023 年 6 条智慧大棚第一季种植记录 & 参数表、2023 基线 & 否则会漏计 12270 斤产量和 7080 元成本 & 已确认 \\
5 & 销售价格只给区间，没有唯一成交价 & 问题 1 以区间中点为基准；上下端用于敏感性；问题 2、3 按题面趋势生成未来价格 & 收益计算、随机情景 & 中点是透明基准，但不是观测成交价 & 已确认 \\
6 & 附件没有 2023 年销售量字段 & 以 2023 年种植面积乘对应亩产量得到作物—季次产量，并将其作为预期销量代理；另做 90\%、100\%、110\% 扫描 & 三问销量基线 & 题面未给销量，不能虚构独立销量数据 & 已确认 \\
7 & 题面中的种植规则分散在文字说明中 & 生成适种矩阵、豆类集合、水浇地两种模式、设施季次集合和 2023 历史启用状态 & 模型可行域 & 将自然语言规则转为可审计的集合与约束 & 已确认 \\
8 & 结果模板包含固定工作表、样式和合并结构 & 问题 1 精确映射目标单元格；问题 2 使用 OOXML 定点回填，禁止覆盖未知结构 & 结果导出 & 保持官方模板结构不变 & 已确认 \\
\bottomrule
\end{longtable}

\section{已确认规则}

\subsection{规则 DC-01：销量代理构造}

设 2023 年地块 \texttt{j}、作物 \texttt{i}、季次 \texttt{s} 的种植面积为 \texttt{x̄\_jis\textasciicircum{}2023}，亩产量为 \texttt{q\_jis\textasciicircum{}2023}，则基准销量代理为

\[
D_{is}^{2023}=\sum_j q_{jis}^{2023}\bar x_{jis}^{2023}.
\]

\noindent 本节汇总见表~\ref{tab:thought-3}。
\begin{longtable}{@{}>{\raggedright\arraybackslash}p{0.26\linewidth}>{\raggedright\arraybackslash}p{0.68\linewidth}@{}}
\caption{思路表第 3 表}\label{tab:thought-3}\\
\toprule
\textbf{口径} & \textbf{处理方式} \\
\midrule
\endfirsthead
\toprule
\textbf{口径} & \textbf{处理方式} \\
\midrule
\endhead
主口径 & 按作物—季次分别形成销量池，与题面“每季销售”表述一致 \\
替代口径 & 按作物汇总全年产量，用于解释敏感性 \\
问题 1 & 2024—2030 年沿用基准销量 \\
问题 2、3 & 按题面增长或波动区间生成各年情景销量 \\
\bottomrule
\end{longtable}

执行要求：必须在正文中称其为“销量代理/预期销量估计”，不能称为附件中的 2023 年真实销售量。

\subsection{规则 DC-02：销售价格与滞销处理}

\noindent 本节汇总见表~\ref{tab:thought-4}。
\begin{longtable}{@{}>{\raggedright\arraybackslash}p{0.17\linewidth}>{\raggedright\arraybackslash}p{0.27\linewidth}>{\raggedright\arraybackslash}p{0.27\linewidth}>{\raggedright\arraybackslash}p{0.23\linewidth}@{}}
\caption{思路表第 4 表}\label{tab:thought-4}\\
\toprule
\textbf{情形} & \textbf{正常销售量} & \textbf{超产部分} & \textbf{收益口径} \\
\midrule
\endfirsthead
\toprule
\textbf{情形} & \textbf{正常销售量} & \textbf{超产部分} & \textbf{收益口径} \\
\midrule
\endhead
问题 1 情形 1 & \texttt{u=min(Q,D)} & 无收入，视为浪费 & \texttt{p·u} \\
问题 1 情形 2 & \texttt{u=min(Q,D)} & 按当年正常售价 50\% 出售 & \texttt{p·u+0.5p(Q-u)} \\
问题 2 主口径 & \(u^\omega{}=min(Q^\omega{},D^\omega{})\) & 无收入，保守处理 & \(p^\omega{}·u^\omega{}\) \\
问题 2 敏感性 & 同上 & 复验 50\% 出售 & 与问题 1 情形 2 同构 \\
问题 3 & 沿用问题 2 所选口径 & 不自行新增清理成本或折价率 & 相关价格和弹性需求下复算 \\
\bottomrule
\end{longtable}

50\% 标准只来源于问题 1 的第二种题设情形，不能用于问题 1 情形 1，也不能未经说明自动推广为问题 2、3 主口径。

\subsection{规则 DC-03：地块与季次结构}

\noindent 本节汇总见表~\ref{tab:thought-5}。
\begin{longtable}{@{}>{\raggedright\arraybackslash}p{0.17\linewidth}>{\raggedright\arraybackslash}p{0.27\linewidth}>{\raggedright\arraybackslash}p{0.27\linewidth}>{\raggedright\arraybackslash}p{0.23\linewidth}@{}}
\caption{思路表第 5 表}\label{tab:thought-5}\\
\toprule
\textbf{地块类型} & \textbf{数量} & \textbf{面积合计} & \textbf{允许结构} \\
\midrule
\endfirsthead
\toprule
\textbf{地块类型} & \textbf{数量} & \textbf{面积合计} & \textbf{允许结构} \\
\midrule
\endhead
平旱地 & 6 & 365 亩 & 每年一季，粮食作物且排除水稻 \\
梯田 & 14 & 619 亩 & 每年一季，粮食作物且排除水稻 \\
山坡地 & 6 & 108 亩 & 每年一季，粮食作物且排除水稻 \\
水浇地 & 8 & 109 亩 & 单季水稻，或第一季普通蔬菜加第二季一种根菜 \\
普通大棚 & 16 & 9.6 亩 & 第一季普通蔬菜，第二季食用菌 \\
智慧大棚 & 4 & 2.4 亩 & 两季普通蔬菜 \\
\bottomrule
\end{longtable}

水浇地第二季根菜集合固定为大白菜、白萝卜、红萝卜；智慧大棚两季均排除这三种根菜。

\subsection{规则 DC-04：轮作与管理约束}

\noindent 本节汇总见表~\ref{tab:thought-6}。
\begin{longtable}{@{}>{\raggedright\arraybackslash}p{0.26\linewidth}>{\raggedright\arraybackslash}p{0.68\linewidth}@{}}
\caption{思路表第 6 表}\label{tab:thought-6}\\
\toprule
\textbf{规则} & \textbf{数学处理} \\
\midrule
\endfirsthead
\toprule
\textbf{规则} & \textbf{数学处理} \\
\midrule
\endhead
不适种组合 & 不创建变量或强制面积为 0 \\
面积使用 & 每个要求种植的地块—年—季面积和等于地块面积 \\
面积—启用联动 & \(\alpha{}_{j} y_{jits} \leq{} x_{jits} \leq{} A_{j} e_{jis} y_{jits}\) \\
最小种植面积 & 主值 \(\alpha{}_{j}=0.5A_{j}\)；扫描 \(\eta{}\in{}{0.25,0.50,0.75,1.00}\) \\
禁止重茬 & 同一地块相邻有效季次不得种同一作物，并纳入 2023 历史边界 \\
三年豆类 & 每个 2023—2030 滚动三年窗口的豆类累计面积不少于地块面积 \\
管理便利性 & 在主收益目标容差内最小化启用次数 \(\sum{}y\) \\
\bottomrule
\end{longtable}

\subsection{规则 DC-05：随机参数身份}

\noindent 本节汇总见表~\ref{tab:thought-7}。
\begin{longtable}{@{}>{\raggedright\arraybackslash}p{0.20\linewidth}>{\raggedright\arraybackslash}p{0.36\linewidth}>{\raggedright\arraybackslash}p{0.38\linewidth}@{}}
\caption{思路表第 7 表}\label{tab:thought-7}\\
\toprule
\textbf{参数} & \textbf{题面区间或趋势} & \textbf{模型处理} \\
\midrule
\endfirsthead
\toprule
\textbf{参数} & \textbf{题面区间或趋势} & \textbf{模型处理} \\
\midrule
\endhead
小麦、玉米销量 & 年增长 5\%—10\% & 按年递推 \\
其他作物销量 & 相对 2023 年约 \ensuremath{\pm{}}5\% & 各年非累积扰动 \\
亩产量 & 年变化 \ensuremath{\pm{}}10\% & 各年非累积扰动 \\
种植成本 & 年增长约 5\% & 主值 5\%，敏感性 4\%—6\% \\
粮食售价 & 基本稳定 & 问题 2 固定；问题 3 加 \ensuremath{\pm{}}1\% 模拟扰动 \\
蔬菜售价 & 年增长约 5\% & 主值 5\%，敏感性 4\%—6\% \\
普通食用菌售价 & 年下降 1\%—5\% & 区间情景递推 \\
羊肚菌售价 & 年下降 5\% & 确定递推 \\
\bottomrule
\end{longtable}

问题 2 的均匀分布和参数独立性、问题 3 的相关载荷、弹性与增产率均是模拟设定，不是附件估计值。

\section{三、第一问解答}

\subsection{1. 数据特征分析}

\noindent 本节汇总见表~\ref{tab:thought-8}。
\begin{longtable}{@{}>{\raggedright\arraybackslash}p{0.20\linewidth}>{\raggedright\arraybackslash}p{0.36\linewidth}>{\raggedright\arraybackslash}p{0.38\linewidth}@{}}
\caption{思路表第 8 表}\label{tab:thought-8}\\
\toprule
\textbf{项目} & \textbf{数据特征} & \textbf{对建模的影响} \\
\midrule
\endfirsthead
\toprule
\textbf{项目} & \textbf{数据特征} & \textbf{对建模的影响} \\
\midrule
\endhead
空间对象 & 54 个地块/大棚，总面积 1213 亩 & 决策必须落实到具体地块，不能只优化作物总面积 \\
作物集合 & 41 种作物，适种类型和季次差异明显 & 需要地块—作物—季次适种矩阵进行变量预筛 \\
历史记录 & 2023 年 87 条种植记录、107 条统计记录 & 同时用于轮作边界、销量代理和参数校验 \\
时间范围 & 2024—2030 共 7 年，包含单季和双季地块 & 需要跨年、跨季联合优化，逐年贪心会破坏轮作可行性 \\
销量限制 & 正常价格销量受 \texttt{D\_is} 限制 & 作物收益跨地块耦合，不能按单亩利润独立排序 \\
管理要求 & 同一作物种植不能过小、过分散 & 需要连续面积变量和 0-1 启用变量联动 \\
\bottomrule
\end{longtable}

第一问的核心任务是求满足全部农业硬约束的七年种植组合，并分别比较两种超产处置规则。

\subsection{2. 数据预处理方案}

\noindent 本节汇总见表~\ref{tab:thought-9}。
\begin{longtable}{@{}>{\raggedright\arraybackslash}p{0.20\linewidth}>{\raggedright\arraybackslash}p{0.36\linewidth}>{\raggedright\arraybackslash}p{0.38\linewidth}@{}}
\caption{思路表第 9 表}\label{tab:thought-9}\\
\toprule
\textbf{处理步骤} & \textbf{具体做法} & \textbf{目的} \\
\midrule
\endfirsthead
\toprule
\textbf{处理步骤} & \textbf{具体做法} & \textbf{目的} \\
\midrule
\endhead
输入身份核验 & 固定附件 SHA-256，核对 54 个对象、41 种作物、1213 亩 & 防止外部保存导致输入版本变化 \\
文本标准化 & 清理空格、换行编码和合并区域空白 & 保证地类、作物和季次键唯一 \\
参数继承 & 补齐智慧大棚第一季蔬菜参数 & 避免少计产量和成本 \\
价格处理 & 区间中点作主价，上下端用于敏感性 & 形成唯一可计算基准 \\
销量构造 & 由 2023 年面积乘亩产量汇总为 \texttt{D\_is} & 弥补附件无销量字段 \\
可行组合预筛 & 只保留 \texttt{e\_jis=1} 的面积和启用变量 & 降低 MILP 规模并消除不适种解 \\
历史边界 & 将 2023 年种植固定为常量 & 正确检查 2024 重茬与三年豆类窗口 \\
\bottomrule
\end{longtable}

\subsection{3. 模型选取及缘由}

主模型选用全周期混合整数线性规划（MILP）与分支定界/分支割算法：

\begin{enumerate}
\item 收益、销量截断、面积、适种、模式互斥和轮作均可线性表达。
\item 连续变量描述种植亩数，二元变量表达是否启用及水浇地模式。
\item 七年一次性优化可避免逐年最优导致后续轮作不可行。
\item 求解器能给出 incumbent、best bound 和 MIP gap，便于区分可行解与已证明最优解。
\item 管理便利性采用字典序次目标，避免用任意权重混合“元”和“启用次数”。
\end{enumerate}

\subsection{4. 模型及其参数介绍}

总产量与正常销量为

\[
Q_{its}=\sum_j q_{jis}x_{jits},\qquad
0\le u_{its}\le Q_{its},\quad u_{its}\le D_{is}.
\]

两种目标分别为

\[
\max Z_1=\sum_{t,i,s}p_{is}u_{its}-\sum_{j,i,t,s}c_{jis}x_{jits},
\]

\[
\max Z_2=\sum_{t,i,s}\left[\tfrac12p_{is}Q_{its}+\tfrac12p_{is}u_{its}\right]
-\sum_{j,i,t,s}c_{jis}x_{jits}.
\]

\noindent 本节汇总见表~\ref{tab:thought-10}。
\begin{longtable}{@{}>{\raggedright\arraybackslash}p{0.20\linewidth}>{\raggedright\arraybackslash}p{0.36\linewidth}>{\raggedright\arraybackslash}p{0.38\linewidth}@{}}
\caption{思路表第 10 表}\label{tab:thought-10}\\
\toprule
\textbf{参数/变量} & \textbf{含义} & \textbf{单位} \\
\midrule
\endfirsthead
\toprule
\textbf{参数/变量} & \textbf{含义} & \textbf{单位} \\
\midrule
\endhead
\texttt{A\_j} & 地块面积 & 亩 \\
\texttt{x\_jits} & 地块—作物—年份—季次种植面积 & 亩 \\
\texttt{y\_jits} & 是否启用该种植组合 & 0-1 \\
\texttt{r\_jt} & 水浇地是否采用单季水稻模式 & 0-1 \\
\texttt{q\_jis} & 亩产量 & 斤/亩 \\
\texttt{c\_jis} & 亩成本 & 元/亩 \\
\texttt{p\_is} & 销售价格 & 元/斤 \\
\texttt{D\_is} & 预期销量代理 & 斤 \\
\texttt{u\_its} & 正常价格销量 & 斤 \\
\bottomrule
\end{longtable}

\subsection{5. 模型训练}

本题不存在机器学习意义上的训练集和梯度训练；“训练”对应 MILP 的构造与求解。

\begin{enumerate}
\item 生成适种变量索引和全部农业硬约束。
\item 对情形 1 最大化 \texttt{Z\_1}，由 HiGHS 执行分支定界/割平面搜索。
\item 对情形 2 仅替换超产收益表达式，保持相同可行域。
\item 在主目标容差内尝试最小化 \(\sum{}y\)；本次正式快照两种情形均未得到第二阶段字典序可行解，因此实际导出的是主阶段 incumbent，而不是已完成的集中化方案。
\item 每个求解点记录 incumbent、best bound、MIP gap、节点数、时间和状态。
\item 固定方案独立复算产量、销量、收入、成本和利润，再写入两个结果模板。
\end{enumerate}

\subsection{6. 模型评分方式}

\noindent 本节汇总见表~\ref{tab:thought-11}。
\begin{longtable}{@{}>{\raggedright\arraybackslash}p{0.20\linewidth}>{\raggedright\arraybackslash}p{0.36\linewidth}>{\raggedright\arraybackslash}p{0.38\linewidth}@{}}
\caption{思路表第 11 表}\label{tab:thought-11}\\
\toprule
\textbf{评价方法} & \textbf{作用} & \textbf{判定方式} \\
\midrule
\endfirsthead
\toprule
\textbf{评价方法} & \textbf{作用} & \textbf{判定方式} \\
\midrule
\endhead
七年总利润 & 衡量经济目标 & 同一滞销口径下越高越好 \\
MIP gap & 衡量当前可行解距可证明界的相对差距 & 达到声明阈值才能称已认证 \\
可行性审计 & 检查农业规则 & 面积、适种、重茬、豆类、模式和整数性均通过 \\
利润独立复算 & 检查模型目标实现 & 复算差不超过数值容差 \\
Excel 回读 & 检查导出正确性 & 模板回读方案与内存方案一致 \\
启用次数 & 衡量种植分散程度 & 在近等利润条件下越少越便于管理 \\
敏感性 & 检查销量、价格和最小面积假设 & 小幅变化不应导致无法解释的结构翻转 \\
\bottomrule
\end{longtable}

\subsection{7. 数值结果分析}

当前自动报告保留的正式求解快照为：

\noindent 本节汇总见表~\ref{tab:thought-12}。
\begin{longtable}{@{}>{\raggedright\arraybackslash}p{0.05\linewidth}>{\raggedright\arraybackslash}p{0.15\linewidth}>{\raggedright\arraybackslash}p{0.27\linewidth}>{\raggedright\arraybackslash}p{0.16\linewidth}>{\raggedright\arraybackslash}p{0.25\linewidth}>{\raggedright\arraybackslash}p{0.08\linewidth}@{}}
\caption{思路表第 12 表}\label{tab:thought-12}\\
\toprule
\textbf{情形} & \textbf{七年利润 incumbent} & \textbf{利润上界} & \textbf{MIP gap} & \textbf{农业约束可行} & \textbf{最优性认证} \\
\midrule
\endfirsthead
\toprule
\textbf{情形} & \textbf{七年利润 incumbent} & \textbf{利润上界} & \textbf{MIP gap} & \textbf{农业约束可行} & \textbf{最优性认证} \\
\midrule
\endhead
超产浪费 & 38,838,192.50 元 & 41,151,153.98 元 & 5.96\% & 是 & 否 \\
超产半价出售 & 63,778,880.31 元 & 63,924,530.53 元 & 0.23\% & 是 & 否 \\
\bottomrule
\end{longtable}

两种方案的正式快照均通过面积、适种、重茬、豆类、水浇地、产量、销量和利润复算审计，但求解状态均为限时退出，且 \texttt{lex\_feasible=False}，不能写成理论全局最优或已完成管理集中化的字典序方案。当前 \texttt{q1\_test/outputs/q1} 中部分工作簿和 \texttt{audit.csv} 后来被短时测试覆盖，其哈希与正式 \texttt{repro.json} 不一致；论文引用前必须重新正式运行并冻结同一批报告、工作簿、审计和复现清单。

\subsection{8. 结论解释与边界}

第一问可以确认：两类滞销处置都能在完整农业规则下形成七年可执行方案；允许超产半价出售显著提高收入，并改变作物组合对销量上限的敏感程度。但表中利润是高质量可行解的 incumbent，不是已经证明达到全局最优的目标值。

销量由 2023 年产量代理，价格采用区间中点，且未计仓储、运输、销毁与换茬成本，因此结论只适用于题面给定的简化经营口径。两种处置情形必须分别解释，不能把 50\% 折价写成所有滞销产品的通用现实规则。

\subsection{9. 补充}

\begin{enumerate}
\item 正式写作前重新运行问题 1，并保证 \texttt{repro.json} 中的输出哈希与当前工作簿、审计和报告一致。
\item 若限时内 gap 仍较大，优先采用更强求解器、MIP start、对称性削减或延长时限，不得用“最优”替代“当前最好可行”。
\item 主结果使用 \(\eta{}=0.5\)；应补充 \(\eta{}=0.25/0.75/1.0\)、销量 90\%/110\% 与价格端点复验。
\item 情形 2 的最优利润理论上不低于情形 1；固定同一方案时也必须满足情形 2 收入不低于情形 1。
\end{enumerate}

\section{四、第二问解答}

\subsection{1. 数据特征分析}

本部分只描述数据现象如何影响“不确定条件下给出最优种植方案”这一目标。

\noindent 本节汇总见表~\ref{tab:thought-13}。
\begin{longtable}{@{}>{\raggedright\arraybackslash}p{0.26\linewidth}>{\raggedright\arraybackslash}p{0.68\linewidth}@{}}
\caption{思路表第 13 表}\label{tab:thought-13}\\
\toprule
\textbf{数据现象} & \textbf{对问题 2 的影响} \\
\midrule
\endfirsthead
\toprule
\textbf{数据现象} & \textbf{对问题 2 的影响} \\
\midrule
\endhead
题面只给未来变化区间或“约”变化，没有概率分布 & 必须区分题面范围与人为分布假设 \\
销量、亩产、成本和价格同时变化 & 单一确定参数方案可能在不利组合下失效 \\
七年共同方案必须事先确定 & \texttt{x,y,r} 不能带情景下标，否则得到的是不可提交的事后方案 \\
农业硬约束数量大 & 每增加一个情景会复制产量、销量和利润变量，计算规模迅速增大 \\
低概率不利情景影响经营稳定 & 只最大化期望利润可能牺牲最差 10\% 情景表现 \\
情景缩减可能删除尾部 & 需要显式保护最低利润层并重新赋权 \\
\bottomrule
\end{longtable}

\subsection{2. 问题分析}

\noindent 本节汇总见表~\ref{tab:thought-14}。
\begin{longtable}{@{}>{\raggedright\arraybackslash}p{0.26\linewidth}>{\raggedright\arraybackslash}p{0.68\linewidth}@{}}
\caption{思路表第 14 表}\label{tab:thought-14}\\
\toprule
\textbf{大类问题} & \textbf{需要回答的具体内容} \\
\midrule
\endfirsthead
\toprule
\textbf{大类问题} & \textbf{需要回答的具体内容} \\
\midrule
\endhead
边际情景生成 & 如何把题面区间转换为可复现的未来参数路径 \\
情景缩减 & 如何从 1000 个原始情景压缩到可求解规模且保留下尾 \\
风险偏好 & 如何同时衡量期望利润与最差 10\% 情景平均利润 \\
共同种植决策 & 如何保证所有情景共享同一套种植面积与启用决策 \\
唯一方案选择 & 如何从 11 个风险权重点中按确定规则选出一套方案 \\
泛化评估 & 如何使用未参与优化的情景评价固定方案 \\
最优性边界 & 如何处理限时有可行解但 gap 未达标的情况 \\
\bottomrule
\end{longtable}

\subsection{3. 指标定义}

\subsubsection{3.1 维持定义}

问题 1 的面积可行性、适种性、重茬、滚动三年豆类、水浇地模式、最小种植面积、利润复算、Excel 回读和 MIP gap 定义全部保留。

\subsubsection{3.2 新增定义}

\noindent 本节汇总见表~\ref{tab:thought-15}。
\begin{longtable}{@{}>{\raggedright\arraybackslash}p{0.20\linewidth}>{\raggedright\arraybackslash}p{0.36\linewidth}>{\raggedright\arraybackslash}p{0.38\linewidth}@{}}
\caption{思路表第 15 表}\label{tab:thought-15}\\
\toprule
\textbf{指标} & \textbf{定义} & \textbf{作用} \\
\midrule
\endfirsthead
\toprule
\textbf{指标} & \textbf{定义} & \textbf{作用} \\
\midrule
\endhead
期望利润 & \(E[\Pi{}]=\sum{}_\omega{} \pi{}_\omega{}\Pi{}_\omega{}\) & 衡量平均经营收益 \\
下尾 CVaR & 利润最低 10\% 情景的加权平均水平 & 衡量不利情景下的收益稳定性 \\
风险权重 \(\lambda{}\) & \(0\leq{}\lambda{}\leq{}1\) & 控制期望利润与下尾利润的偏好 \\
风险前沿 & 各 \(\lambda{}\) 对应的期望利润—下尾利润组合 & 展示风险收益取舍 \\
样本外均值/标准差 & 5000 个新情景下固定方案统计量 & 衡量未参与优化情景中的平均与波动 \\
10\% 分位数 & 利润经验分布的第 10 百分位 & 描述较差情景阈值 \\
最差 10\% 平均利润 & 样本外最低 10\% 利润均值 & 与模型 LCVaR 对照 \\
亏损概率 & \texttt{P(\ensuremath{\Pi{}}<0)} & 衡量出现负利润的频率 \\
\bottomrule
\end{longtable}

\subsubsection{3.3 变更定义}

问题 1 的单一确定利润目标变更为

\[
\max Z_\lambda=(1-\lambda)\sum_\omega\pi_\omega\Pi_\omega
+\lambda\left[\zeta-\frac{1}{1-\beta}\sum_\omega\pi_\omega\xi_\omega\right],
\]

\[
\xi_\omega\ge \zeta-\Pi_\omega,\qquad \xi_\omega\ge0,
\]

其中 \(\beta{}=0.90\)。\(\lambda{}=0\) 为风险中性，\(\lambda{}=1\) 只关注下尾利润。超产浪费是问题 2 主口径，半价出售仅作敏感性。

\subsection{4. 模型与算法组件及选取缘由}

\noindent 本节汇总见表~\ref{tab:thought-16}。
\begin{longtable}{@{}>{\raggedright\arraybackslash}p{0.17\linewidth}>{\raggedright\arraybackslash}p{0.27\linewidth}>{\raggedright\arraybackslash}p{0.27\linewidth}>{\raggedright\arraybackslash}p{0.23\linewidth}@{}}
\caption{思路表第 16 表}\label{tab:thought-16}\\
\toprule
\textbf{模型或算法} & \textbf{是什么} & \textbf{选取缘由} & \textbf{解决的问题} \\
\midrule
\endfirsthead
\toprule
\textbf{模型或算法} & \textbf{是什么} & \textbf{选取缘由} & \textbf{解决的问题} \\
\midrule
\endhead
Latin Hypercube Sampling & 每一维在 \texttt{[0,1]} 等概率分层后映射到题面范围 & 比普通随机抽样更均匀覆盖各边际区间 & 原始情景生成 \\
利润分层 PAM k-medoids & 先按 Q1 基线代理利润分十层，再用 L1 距离选真实 medoid & 保留真实样本并保护最低利润层 & 情景缩减 \\
均值—下尾 CVaR 随机 MILP & 在共同农业决策上优化期望与尾部利润 & CVaR 可线性化并继承 Q1 可行域 & 风险优化 \\
三级字典序 & 风险目标、期望利润、最少启用依次优化 & 消除纯尾部目标退化并保证管理可解释 & 唯一可执行方案 \\
膝点/99\%规则 & 按归一化风险前沿的几何结构自动选择 & 避免人工选择最好看的 \(\lambda{}\) & 唯一方案选择 \\
样本外 Monte Carlo & 固定方案在新随机流中向量化复算 & 区分优化表现和泛化表现 & 稳健性验证 \\
\bottomrule
\end{longtable}

\subsection{5. 整体路线/算法结构}

\begin{center}
\begin{minipage}{0.78\linewidth}
\small\raggedright
问题1清洗数据与确定性农业可行域\par
        \ensuremath{\downarrow} 题面区间映射\par
1000个LHS原始独立参数情景\par
        \ensuremath{\downarrow} Q1固定方案代理利润分层\par
尾部保护PAM k-medoids缩减为30个代表情景\par
        \ensuremath{\downarrow} \ensuremath{\lambda{}}=0,0.1,...,1\par
均值—下尾CVaR随机MILP\par
        \ensuremath{\downarrow} 每个\ensuremath{\lambda{}}执行三级字典序\par
风险前沿与唯一方案选择\par
        \ensuremath{\downarrow} 5000个未参与优化的新情景\par
样本外均值、尾部、最低利润和亏损概率\par
        \ensuremath{\downarrow}\par
农业约束、利润、CVaR、模板与复现审计\par
\end{minipage}
\end{center}

\subsection{6. 训练方案}

\subsubsection{6.1 第一次训练：边际情景生成}

使用种子 2024 生成 1000 个 LHS 原始情景。小麦、玉米销量按 5\%—10\% 年增长递推；其他销量和亩产分别按 \ensuremath{\pm{}}5\%、\ensuremath{\pm{}}10\% 非累积扰动；成本、蔬菜价格和菌类价格按题面趋势递推。问题 2 基准暂按参数相互独立、区间内均匀分布，必须在报告中标注为模拟假设。

\subsubsection{6.2 第二次训练：情景缩减与风险网格}

用问题 1 超产浪费方案复算每个原始情景的代理利润，按利润分成十个等频层，最低 10\% 层至少保留 \texttt{ceil(0.1K)} 个代表情景。各层内采用 PAM k-medoids 和 L1 距离缩减，代表情景权重为所属簇样本比例。随后对 \(\lambda{}=0,0.1,...,1\) 的 11 个点分别求解。

\subsubsection{6.3 第三次训练：三级字典序与唯一方案}

每个 \(\lambda{}\) 依次最大化风险目标、在金额容差内最大化期望利润、再最小化启用次数。存在完整前沿时，最大归一化垂距至少 0.02 则选膝点；否则选下尾利润达到前沿最大值 99\% 的最小 \(\lambda{}\)，再按期望利润和启用次数破同分。没有有限 incumbent 的点不得进入前沿。

\subsection{7. 验证方案}

\begin{enumerate}
\item 对每个代表情景独立复算 \(Q^\omega{},u^\omega{},\Pi{}_\omega{}\) 和 LCVaR。
\item 检查共同决策不带情景下标，全部农业硬约束最大违约不超过数值容差。
\item 优化与样本外使用不同随机流；样本外 5000 个情景只复算，不重新优化。
\item 比较 \texttt{K=20/30/50}；关键指标变化超过 2\% 时扩大代表情景数。
\item 用区间端点联合压力情景、三角分布和至少 5 个随机种子检验稳定性。
\item 报告 incumbent、bound、MIP gap、时间和求解状态；前沿不完整时禁止声称已完成全风险优化。
\item OOXML 回填后检查非目标结构不变，回读面积与内存方案一致。
\end{enumerate}

\subsection{8. 结论对比}

\noindent 本节汇总见表~\ref{tab:thought-17}。
\begin{longtable}{@{}>{\raggedright\arraybackslash}p{0.05\linewidth}>{\raggedright\arraybackslash}p{0.15\linewidth}>{\raggedright\arraybackslash}p{0.27\linewidth}>{\raggedright\arraybackslash}p{0.16\linewidth}>{\raggedright\arraybackslash}p{0.25\linewidth}>{\raggedright\arraybackslash}p{0.08\linewidth}@{}}
\caption{思路表第 17 表}\label{tab:thought-17}\\
\toprule
\textbf{方案/权重点} & \textbf{期望利润} & \textbf{下尾 CVaR} & \textbf{激活数} & \textbf{状态} & \textbf{MIP gap} \\
\midrule
\endfirsthead
\toprule
\textbf{方案/权重点} & \textbf{期望利润} & \textbf{下尾 CVaR} & \textbf{激活数} & \textbf{状态} & \textbf{MIP gap} \\
\midrule
\endhead
\(\lambda{}=0.0\) & 24,376,595 元 & 24,054,552 元 & 574 & 可行但未证明 & 5.23\% \\
\(\lambda{}=0.1\) & 22,973,943 元 & 22,660,189 元 & 574 & 可行但未证明 & 5.23\% \\
\(\lambda{}=0.2...1.0\) & 无有限可行 incumbent & 无 & 0 & 限时未找到可行解 & 100\% \\
\bottomrule
\end{longtable}

当前只得到两个可用点，且 \(\lambda{}=0\) 同时支配 \(\lambda{}=0.1\) 的期望利润与下尾利润；因此程序选择 \(\lambda{}=0\)，但这只是当前不完整前沿中的唯一非支配可用方案，不能解释为风险权衡已经充分搜索。

\subsection{9. 最终结果}

选中方案的 5000 个样本外情景结果为：

\noindent 本节汇总见表~\ref{tab:thought-18}。
\begin{longtable}{@{}>{\raggedright\arraybackslash}p{0.26\linewidth}>{\raggedright\arraybackslash}p{0.68\linewidth}@{}}
\caption{思路表第 18 表}\label{tab:thought-18}\\
\toprule
\textbf{指标} & \textbf{数值} \\
\midrule
\endfirsthead
\toprule
\textbf{指标} & \textbf{数值} \\
\midrule
\endhead
平均利润 & 24,377,039 元 \\
标准差 & 118,385 元 \\
10\% 分位数 & 24,227,417 元 \\
最差 10\% 平均利润 & 24,172,747 元 \\
最低利润 & 23,964,067 元 \\
亏损概率 & 0 \\
\bottomrule
\end{longtable}

农业约束、利润复算、CVaR 复算和 OOXML 回读均通过；但 \texttt{frontier\_complete=False}、\texttt{selected\_certified=False}。最终文稿只能称其为“当前求解条件下的可行候选方案”或 \texttt{q2\_feasible\_baseline}。

\subsection{10. 补充}

\begin{enumerate}
\item 当前 11 个 \(\lambda{}\) 中 9 个没有可行解，正式交付前应使用支持 MIP start 的求解器、增加时限或先降低 \texttt{K} 建立完整前沿。
\item \(\lambda{}=0.1\) 的下尾利润低于 \(\lambda{}=0\)，主要反映限时可行解质量，不应据此断言风险厌恶必然降低尾部表现。
\item 均匀分布和独立性不是题面事实；需补充三角分布、端点和多种子敏感性。
\item 问题 2 基线可用于问题 3 公平复算，但标签必须保留“可行、未认证”。
\end{enumerate}

\section{五、第三问解答}

\subsection{1. 数据特征分析}

\noindent 本节汇总见表~\ref{tab:thought-19}。
\begin{longtable}{@{}>{\raggedright\arraybackslash}p{0.26\linewidth}>{\raggedright\arraybackslash}p{0.68\linewidth}@{}}
\caption{思路表第 19 表}\label{tab:thought-19}\\
\toprule
\textbf{数据特征} & \textbf{对问题 3 的影响} \\
\midrule
\endfirsthead
\toprule
\textbf{数据特征} & \textbf{对问题 3 的影响} \\
\midrule
\endhead
附件只有 2023 年单年截面 & 无法从数据可靠估计 41 种作物的相关系数、价格弹性或轮作增产率 \\
销量、价格、成本和亩产存在共同宏观与气候驱动 & 问题 2 的独立抽样可能低估联合不利情景 \\
同类作物存在替代关系 & 某作物相对价格变化会改变其他作物的需求 \\
豆类根瘤菌有利于后茬作物 & 可在生产侧建立有题面依据的跨年互补增益 \\
农业约束与问题 2 相同 & 可以复用共同决策与随机 MILP 可行域 \\
Q2 基线未认证 & 比较时只能把它作为固定可行候选，不得称为理论最优基准 \\
\bottomrule
\end{longtable}

\subsection{2. 问题分析}

\noindent 本节汇总见表~\ref{tab:thought-20}。
\begin{longtable}{@{}>{\raggedright\arraybackslash}p{0.26\linewidth}>{\raggedright\arraybackslash}p{0.68\linewidth}@{}}
\caption{思路表第 20 表}\label{tab:thought-20}\\
\toprule
\textbf{大类问题} & \textbf{具体任务} \\
\midrule
\endfirsthead
\toprule
\textbf{大类问题} & \textbf{具体任务} \\
\midrule
\endhead
联合相关情景 & 在保留 Q2 各边际区间的同时生成变量间和跨年的相关秩结构 \\
替代性 & 用稀疏交叉价格弹性把价格冲击转化为需求变化 \\
互补性 & 用豆类前茬变量表达下一年度非豆类亩产增益 \\
模型线性性 & 保证交互机制预处理或线性化后仍可由 MILP 求解 \\
参数不可识别 & 明确载荷、弹性和增产率属于模拟设定并做弱中强扫描 \\
公平比较 & Q2 与 Q3 固定方案必须使用同一批相关样本外情景和共同随机数 \\
失稳回退 & 若推荐结构对模拟参数高度敏感，改报条件性策略或转向分布鲁棒模型 \\
\bottomrule
\end{longtable}

\subsection{3. 指标定义}

\subsubsection{3.1 维持定义}

沿用 Q2 的期望利润、下尾 CVaR、风险前沿、样本外利润指标、三级字典序、农业硬约束、利润复算、MIP gap 和模板审计。

\subsubsection{3.2 新增定义}

\noindent 本节汇总见表~\ref{tab:thought-21}。
\begin{longtable}{@{}>{\raggedright\arraybackslash}p{0.20\linewidth}>{\raggedright\arraybackslash}p{0.36\linewidth}>{\raggedright\arraybackslash}p{0.38\linewidth}@{}}
\caption{思路表第 21 表}\label{tab:thought-21}\\
\toprule
\textbf{指标} & \textbf{定义} & \textbf{作用} \\
\midrule
\endfirsthead
\toprule
\textbf{指标} & \textbf{定义} & \textbf{作用} \\
\midrule
\endhead
目标 Kendall \(\tau{}*\) & \(2/\pi{}·arcsin(R_{lat})\) & 审计 t-Copula 的秩相关结构 \\
相关误差 & 审计对上样本 Kendall 与目标值的最大绝对差 & 判断原始及缩减情景是否保留依赖 \\
弹性稳定条件 & \(\sum{}_{h\neq{}i}\lvert e_{ih}\rvert\leq{}\lvert e_{ii}\rvert\) & 防止交叉效应支配自身价格效应 \\
互补激活误差 & \(\lvert w-xb\rvert\) 的最大值 & 检查乘积线性化正确性 \\
成对利润差 & \(\Delta{}\Pi{}_\omega{}=\Pi{}_\omega{}^Q3-\Pi{}_\omega{}^Q2\) & 在相同情景中比较两套方案 \\
面积 L1 距离 & 两方案种植面积绝对差之和 & 衡量策略结构变化 \\
激活 Jaccard & 两方案非零种植组合交并比 & 衡量种植组合稳定性 \\
参数稳定频率 & 弱中强档中主要结构保持一致的比例 & 判断能否给单一推荐 \\
\bottomrule
\end{longtable}

\subsubsection{3.3 变更定义}

Q2 的独立边际情景改为相关联合情景；基础需求经交叉弹性修正：

\[
D_{its}^{\omega}=D_{its}^{base,\omega}
\exp\left(\sum_{h\in I_s}e_{ih}^{(s)}
\ln\frac{p_{hts}^{\omega}}{\bar p_{hts}}\right).
\]

豆类前茬增益后的产量为

\[
Q_{its}^{\omega}=\sum_jq_{jits}^{\omega}
\left(x_{jits}+\gamma_iw_{jits}\right),\qquad w_{jits}=x_{jits}b_{jt}.
\]

最终仍优化 Q2 的均值—下尾 CVaR 目标，但 \(\Pi{}_\omega{}\) 使用相关需求、价格、成本、亩产和互补增益重新计算。

\subsection{4. 模型与算法组件及选取缘由}

\noindent 本节汇总见表~\ref{tab:thought-22}。
\begin{longtable}{@{}>{\raggedright\arraybackslash}p{0.17\linewidth}>{\raggedright\arraybackslash}p{0.27\linewidth}>{\raggedright\arraybackslash}p{0.27\linewidth}>{\raggedright\arraybackslash}p{0.23\linewidth}@{}}
\caption{思路表第 22 表}\label{tab:thought-22}\\
\toprule
\textbf{模型或算法} & \textbf{是什么} & \textbf{选取缘由} & \textbf{解决的问题} \\
\midrule
\endfirsthead
\toprule
\textbf{模型或算法} & \textbf{是什么} & \textbf{选取缘由} & \textbf{解决的问题} \\
\midrule
\endhead
因子 t-Copula & 宏观、气候和类别因子与共享卡方尺度共同产生厚尾相关 & 能以少量透明因子构造半正定联合结构 & 参数及跨年相关 \\
秩保持 LHS 重排 & 按相关 t 分数秩次重排各维 LHS 边际 & 保留 Q2 边际分层与题面范围 & 边际—相关兼容 \\
稀疏交叉价格弹性 & 同季同类少量弹性边，未获依据者为 0 & 避免 41 维稠密矩阵不可识别 & 作物替代性 \\
豆类前茬线性化 & 用 \texttt{b} 与 \texttt{w} 表示是否有豆类前茬及受益面积 & 直接利用题面互补机制且保持线性 & 生产互补性 \\
相关/尾部保护缩减 & 距离同时包含参数、因子和代理利润 & 避免缩减后丢失联合尾部和相关方向 & 计算规模控制 \\
均值—下尾 CVaR 随机 MILP & 复用 Q2 风险模型并替换情景生成器 & 模型接口清晰且可与 Q2 公平比较 & 最终种植优化 \\
\bottomrule
\end{longtable}

\subsection{5. 整体路线/算法结构}

\begin{center}
\begin{minipage}{0.78\linewidth}
\small\raggedright
Q2题面边际与共同农业可行域\par
        \ensuremath{\downarrow} 稳态AR(1)因子 + 每情景共享卡方尺度\par
因子t-Copula七年相关秩分数\par
        \ensuremath{\downarrow} 按秩重排各维LHS样本\par
保留边际范围的相关联合情景\par
        \ensuremath{\downarrow} 交叉价格弹性\par
情景最终需求\par
        \ensuremath{\downarrow} 豆类前茬b与受益面积w线性化\par
相关/尾部保护情景缩减\par
        \ensuremath{\downarrow} 均值—下尾CVaR三级字典序MILP\par
Q3风险前沿与唯一候选方案\par
        \ensuremath{\downarrow} Q2固定方案与Q3方案共同随机数复算\par
四组消融、成对Bootstrap与弱中强敏感性\par
\end{minipage}
\end{center}

\subsection{6. 训练方案}

\subsubsection{6.1 多因素相关情景训练}

因子初始状态固定为 \texttt{g\_k,2024\textasciitilde{}N(0,1)}，2025—2030 按

\[
g_{k,t}=\rho_kg_{k,t-1}+\sqrt{1-\rho_k^2}\nu_{k,t}
\]

递推。每个情景只抽取一个 \(S_\omega{}~\chi{}^{2}(ν)\)，由该情景内全部变量和七个年份共享。基准取 \texttt{ν=5}、载荷强度 1、\(\rho{}=0.5\)；相关强度和持续性按弱中强档扫描。审计对固定为所有 \(|\tau{}*|\geq{}0.10\) 的非对角无序对，加上同一维度相邻年份对。

\subsubsection{6.2 替代与互补组件训练}

粮食、蔬菜、食用菌基准自身价格弹性分别为 \texttt{-0.25/-0.50/-0.60}，同类替代弹性行总和分别为 \texttt{0.15/0.20/0.25}，整体乘 \(\kappa{}_{E}\in{}{0.5,1.0,1.5}\)。非豆类作物豆类前茬增产率取 \(\gamma{}\in{}{0,0.03,0.06}\)，豆类自身为 0。上述参数只用于机制模拟，不能写成附件回归估计。

\subsubsection{6.3 联合随机规划与消融训练}

基准建议生成 2000 个相关原始情景，比较 \texttt{K=20/30/50}，并对 \(\lambda{}=0...1\) 运行风险前沿。四组消融为：Q2 冻结方案只复算、仅相关性、仅替代/互补、完整 Q3。后三组保持 \(\beta{}\)、\(\lambda{}\) 网格、场景规模、时限、gap 和选择规则一致；全部方案使用同一批 Q3 样本外情景。

\subsection{7. 验证方案}

\begin{enumerate}
\item 检查潜在线性相关矩阵对称、对角为 1、最小特征值不低于 \texttt{-1e-10}。
\item 原始样本审计对上 Kendall 最大误差不超过 0.05；缩减样本按情景权重计算加权 Kendall，方向一致且最大误差不超过 0.15。
\item 检查所有边际范围、递推关系、弹性符号和行稳定条件。
\item 用微型算例手算验证 \texttt{w=xb}：\texttt{b=0} 时 \texttt{w=0}，\texttt{b=1} 时 \texttt{w=x}。
\item 农业硬约束、利润、CVaR、整数性和结果导出独立复算。
\item Q2 与 Q3 使用共同随机数，对 \texttt{mean(\ensuremath{\Delta{}}\ensuremath{\Pi{}})} 与下尾差做成对 bootstrap 95\% 区间。
\item Q3 的平均利润损失预设不超过 Q2 的 2\%，并同时报告 1\%、3\%、5\% 门槛敏感性。
\item 若弱中强档下主要作物面积频繁翻转，不给单一无条件推荐，转入分布鲁棒回退评估。
\end{enumerate}

\subsection{8. 结论对比}

\noindent 本节汇总见表~\ref{tab:thought-23}。
\begin{longtable}{@{}>{\raggedright\arraybackslash}p{0.12\linewidth}>{\raggedright\arraybackslash}p{0.20\linewidth}>{\raggedright\arraybackslash}p{0.24\linewidth}>{\raggedright\arraybackslash}p{0.24\linewidth}>{\raggedright\arraybackslash}p{0.14\linewidth}@{}}
\caption{思路表第 23 表}\label{tab:thought-23}\\
\toprule
\textbf{比较组} & \textbf{相关性} & \textbf{替代/互补} & \textbf{是否重新优化} & \textbf{当前结果状态} \\
\midrule
\endfirsthead
\toprule
\textbf{比较组} & \textbf{相关性} & \textbf{替代/互补} & \textbf{是否重新优化} & \textbf{当前结果状态} \\
\midrule
\endhead
Q2 可行候选基线 & 独立边际生成 & 无 & 否，只在 Q3 情景中复算 & 已有固定方案，未认证 \\
仅相关性 & 有 & 无 & 是 & 待实现 \\
仅替代/互补 & 无 & 有 & 是 & 待实现 \\
完整 Q3 & 有 & 有 & 是 & 待实现 \\
\bottomrule
\end{longtable}

预期对比指标包括期望利润、标准差、10\% 分位数、最差 10\% 平均利润、最低利润、亏损概率、面积 L1 距离、激活 Jaccard、启用次数和集中度。目前没有真实 Q3 运行值，不能提前判断哪一组最好。

\subsection{9. 最终结果}

问题 3 当前只完成建模方案、工程结构和编程手实现合同，尚未完成 Python 求解链、正式风险前沿、样本外比较和结果工作簿。因此：

\noindent 本节汇总见表~\ref{tab:thought-24}。
\begin{longtable}{@{}>{\raggedright\arraybackslash}p{0.26\linewidth}>{\raggedright\arraybackslash}p{0.68\linewidth}@{}}
\caption{思路表第 24 表}\label{tab:thought-24}\\
\toprule
\textbf{交付项} & \textbf{当前状态} \\
\midrule
\endfirsthead
\toprule
\textbf{交付项} & \textbf{当前状态} \\
\midrule
\endhead
数学模型 & 已冻结并通过 M1 建模质检 \\
工程骨架 & 已建立 \texttt{q3\_test} 目录与接口规划 \\
求解代码 & 未实现 \\
P1 最小纵向测试 & 未执行 \\
Q3 正式数值解 & 不存在 \\
Q2/Q3 消融和成对比较 & 未执行 \\
可提交结果表 & 不存在 \\
\bottomrule
\end{longtable}

只有在 Q3 方案农业可行、相关审计通过、下尾改善的成对 95\% 区间下界不小于 0、平均利润损失满足预设门槛且求解状态完整披露后，才能给出稳健性推荐。

\subsection{10. 补充}

\begin{enumerate}
\item 下游编程手应先阅读 \texttt{q3\_test/AGENT.md}，按 P1 最小闭环实现，不得先跑全量模型。
\item 求解器需动态探测 MIP start、bound 和检查点能力；能力不足时显式降级并返回相应退出码。
\item 加权 Kendall 的 concordance/discordance 与 ties 处理必须在代码注释和测试中冻结，避免统计库口径差异。
\item 题目没有提供真实弹性或相关估计依据，论文必须把这些数值标为情景模拟参数。
\item 问题 3 没有专用官方模板，候选工作簿只能作为展示文件，不冒充题目指定提交格式。
\end{enumerate}

\section{六、补充验证与最终诊断}

\subsection{1. 问题 1 产物一致性诊断}

\begin{itemize}
\item \texttt{Q1\_建模实现报告.md} 与 \texttt{repro.json} 记录了 600 秒正式求解快照。
\item 当前 \texttt{outputs/q1} 中部分文件在 2026-08-28 14:13—14:14 被后续短时运行覆盖，而正式报告与复现清单生成于 00:21。
\item 因此问题 1 的模型思路和正式报告数值可作阶段性参考，但论文交付前必须重跑并重新冻结同一批文件哈希。
\item 在重新冻结前，不应把当前工作簿、当前 \texttt{audit.csv} 与旧正式报告混合作为同一次运行证据。
\end{itemize}

\subsection{2. 问题 2 风险前沿诊断}

\begin{itemize}
\item 原始 1000、缩减 30、样本外 5000 个情景，\(\beta{}=0.90\)。
\item 只有 \(\lambda{}=0\) 与 \(\lambda{}=0.1\) 得到有限可行方案；其余 9 点限时无可行 incumbent。
\item 当前选择 \(\lambda{}=0\) 的原因是它是可用点中的唯一非支配方案，不代表完整前沿上的理论最佳权重。
\item 选中方案农业约束和样本外复算通过，但 MIP gap 约 5.23\%，\texttt{frontier\_complete=False}、\texttt{selected\_certified=False}。
\end{itemize}

\subsection{3. 问题 3 参数身份与公平比较诊断}

\begin{itemize}
\item t-Copula 载荷、AR(1) 系数、自由度、交叉价格弹性和豆类前茬增产率均属于透明模拟设置。
\item Q2 固定方案不重跑 \(\lambda{}\)，只在同一批 Q3 相关样本外情景中复算；其余三组才重新优化。
\item 使用共同随机数和逐情景利润差，避免把不同随机样本上的目标值直接比较。
\item 若相关、弹性或增产率轻微改变就造成主要种植结构翻转，最终只能给出条件性方案。
\end{itemize}

\subsection{4. 三问统一完成门槛}

\noindent 本节汇总见表~\ref{tab:thought-25}。
\begin{longtable}{@{}>{\raggedright\arraybackslash}p{0.17\linewidth}>{\raggedright\arraybackslash}p{0.27\linewidth}>{\raggedright\arraybackslash}p{0.27\linewidth}>{\raggedright\arraybackslash}p{0.23\linewidth}@{}}
\caption{思路表第 25 表}\label{tab:thought-25}\\
\toprule
\textbf{门槛} & \textbf{问题 1} & \textbf{问题 2} & \textbf{问题 3} \\
\midrule
\endfirsthead
\toprule
\textbf{门槛} & \textbf{问题 1} & \textbf{问题 2} & \textbf{问题 3} \\
\midrule
\endhead
输入哈希与只读附件 & 已建立 & 已建立 & 已规划继承 \\
农业硬约束审计 & 已通过历史正式快照 & 已通过 & 待执行 \\
利润/CVaR独立复算 & 利润已通过历史快照 & 已通过 & 待执行 \\
最优性/前沿认证 & 未通过 & 未通过 & 待执行 \\
样本外评估 & 不适用/敏感性为主 & 已完成 5000 情景 & 待执行 \\
参数敏感性 & 部分完成，需重冻结 & 尚不完整 & 待执行 \\
结果与复现哈希一致 & 当前不一致，需重跑 & 当前基线已冻结 & 待执行 \\
可直接写入最终论文 & 否 & 仅可写“可行未认证” & 否 \\
\bottomrule
\end{longtable}

本表只作为论文写作底稿和工程交接依据。最终论文中的数值、状态和结论必须再次与同一次正式运行的代码输出、审计、图表和复现清单逐项核对。
\end{document}
